\section{Related Work}
\label{sec:related}
\paragraph{Inductive Reasoning with Language Models} In the context of language models, inductive reasoning involves generating rules from observed examples that capture the underlying data-generating mechanism~\citep{wang2023hypothesis, qiu2023phenomenal, yang-etal-2024-language}, with diverse downstream applications in domains such as coding~\citep{shao-etal-2025-case2code} and dialogue systems~\citep{xie-etal-2024-shot-dialogue}. While several recent works benchmark and evaluate language models' inductive reasoning abilities~\citep{li2025mirageevaluatingexplaininginductive, gendron2024largelanguagemodelsstrong}, how to systematically enhance the inductive reasoning capabilities of language models, as studied in this work, remains an open problem. While existing works aim to improve inductive reasoning using frozen models through hypothesis search~\citep{wang2023hypothesis} and refinement~\citep{qiu2023phenomenal}, these methods rely on large proprietary models, which can be prohibitively expensive for inference. We demonstrate that systematically improving the inductive reasoning ability of open-weight models through distillation can match the accuracy of large proprietary models while allowing significantly cheaper inference.

\paragraph{Distillations for Language Model Training} 
Recently, there has been increasing interest in distilling domain-specific or task-specific knowledge and abilities from large language models. This paradigm has been applied to domains such as summarization~\citep{jung-etal-2024-impossible}, commonsense reasoning~\citep{west-etal-2022-symbolic}, chain-of-thought reasoning~\citep{li-etal-2024-mode}, among others. However, it is unclear whether this paradigm can be applied to the distillation of inductive reasoning abilities from language models, especially if the quality of intermediate reasoning steps generated by the model is not evaluated. We address this issue in our work and demonstrate that effective evaluation of the hypothesis-generation step can lead to improved accuracy and reduced cost for the resulting distilled models.

\paragraph{Learning from In-context Demonstrations}
ICL has emerged as a powerful paradigm for language model reasoning~\citep{dong-etal-2024-survey}. A growing body of research highlights the significance of demonstration selection for ICL. Prior work finds that the number, order, and diversity of in-context examples significantly affect model performance~\citep{min-etal-2022-rethinking, levy-etal-2023-diverse, peng-etal-2024-revisiting}. Retrieval has been shown to be an effective method for demonstration selection~\citep{li-etal-2023-unified, DBLP:conf/nips/ZhangZ023, DBLP:journals/corr/abs-2401-11624, gao2025learning}. Additionally, many studies propose methods to enhance the adaptability and generalization of ICL~\citep{metaICL,DBLP:conf/icml/LiIPO23,xu-etal-2023-unleash, wei-etal-2023-symbol}.

